\chapter{Fundamentos teóricos}
% objetivo del laboratorio

El propósito de este laboratorio es la verificación experimental de la condición de equilibrio rotacional para un cuerpo que puede rotar en torno a un eje fijo. Para ello se debe verificar que si el cuerpo no gira, es porque la suma de los torques causados por las fuerzas externas que actúan sobre él es cero.

% que es el torque
El momento de fuerza, a menudo también llamado <<par de torsión>> o <<torque>> y representado por la letra griega <<$\tau$>>, es una magnitud vectorial que indica la tendencia de una fuerza a producir o modificar la rotación de un cuerpo alrededor de un eje.

Vectorialmente, el torque se define como el producto vectorial entre el vector posición $\vec{r}$, medido desde el eje de rotación hasta el punto de aplicación de la fuerza, y el vector fuerza $\vec{F}$.

\begin{equation}
\vec{\tau}=\vec{r}\times\vec{F}
\end{equation}

% formula del torque
La magnitud del torque está dada por:

\begin{equation}
    {\tau} = rF\sin{(\theta)}
\end{equation}

Cabe destacar que se considera un par de torsión positivo cuando el cuerpo gira en sentido anti-horario y un par de torsión negativo cuando el cuerpo gira en sentido horario.

% intro a sumatoria de torque
Un cuerpo puede ser afectado por diversas fuerzas, lo que implica que puede ser influenciado por uno o más momentos de fuerza que pueden alterar su comportamiento en cuanto a su rotación.

% equilibrio rotacional
El equilibrio rotacional consiste en el estado obtenido por un cuerpo cuando todos los momentos de fuerza que afectan al cuerpo se anulan entre sí; matemáticamente, esto se describe como la igualdad entre la sumatoria de torques y cero.

% formula de equilibrio rotacional
Es decir:

\begin{equation}
    \sum \tau = 0
\end{equation}

\newpage

% torque de los pesos
En el contexto del experimento, los torques corresponden a aquellos producidos por el peso de las masas respecto a un eje de rotación. Debido a que la fuerza es perpendicular al brazo de la palanca, el ángulo entre ambos vectores es de $90^\circ$, por lo que $\sin(90^\circ)=1$.

Por lo tanto, el torque producido por cada peso está dado por:

\begin{equation}
    \tau = rmg
\end{equation}